\documentclass[11pt,a4paper]{article}

\usepackage[utf8]{inputenc}
\usepackage[T1]{fontenc}
\usepackage[margin=2.5cm]{geometry}
\usepackage{hyperref}
\usepackage{xcolor}
\usepackage{listings}
\usepackage{enumitem}
\usepackage{booktabs}
\usepackage{longtable}
\usepackage{amsmath,amssymb}
\usepackage{fancyhdr}
\usepackage{titlesec}

\definecolor{isls-blue}{HTML}{1a5276}
\definecolor{isls-orange}{HTML}{e67e22}
\definecolor{isls-green}{HTML}{27ae60}
\definecolor{isls-red}{HTML}{c0392b}
\definecolor{isls-gray}{HTML}{7f8c8d}
\definecolor{code-bg}{HTML}{f7f9fb}

\lstdefinestyle{rust}{
  language=C,
  basicstyle=\ttfamily\small,
  keywordstyle=\color{isls-blue}\bfseries,
  commentstyle=\color{isls-gray}\itshape,
  stringstyle=\color{isls-green},
  backgroundcolor=\color{code-bg},
  frame=single,
  rulecolor=\color{isls-gray!30},
  breaklines=true,
  showstringspaces=false,
  tabsize=4,
  morekeywords={fn,let,mut,pub,struct,impl,enum,use,mod,self,Self,
    crate,super,where,trait,for,in,if,else,match,return,async,await,
    Ok,Err,Some,None,Vec,String,HashMap,HashSet,PathBuf,Result,Option},
}
\lstdefinestyle{bash}{
  language=bash,
  basicstyle=\ttfamily\small,
  backgroundcolor=\color{code-bg},
  frame=single,
  rulecolor=\color{isls-gray!30},
  breaklines=true,
  showstringspaces=false,
}
\lstdefinestyle{toml}{
  basicstyle=\ttfamily\small,
  backgroundcolor=\color{code-bg},
  frame=single,
  rulecolor=\color{isls-gray!30},
  breaklines=true,
  showstringspaces=false,
  commentstyle=\color{isls-gray}\itshape,
  morecomment=[l]{\#},
}

\pagestyle{fancy}
\fancyhf{}
\fancyhead[L]{\small\textcolor{isls-gray}{ISLS D8 Spec --- Offener Horizont}}
\fancyhead[R]{\small\textcolor{isls-gray}{v1.0.0 --- \today}}
\fancyfoot[C]{\thepage}

\titleformat{\section}
  {\Large\bfseries\color{isls-blue}}{\thesection}{1em}{}
\titleformat{\subsection}
  {\large\bfseries\color{isls-blue!80}}{\thesubsection}{1em}{}
\titleformat{\subsubsection}
  {\normalsize\bfseries\color{isls-blue!60}}{\thesubsubsection}{1em}{}

\hypersetup{
  colorlinks=true,
  linkcolor=isls-blue,
  urlcolor=isls-orange,
  citecolor=isls-green,
}

\begin{document}

\begin{titlepage}
\centering
\vspace*{2.5cm}

{\Huge\bfseries\color{isls-blue} ISLS D8 Specification}\\[0.5cm]
{\LARGE\color{isls-orange} Offener Horizont}\\[0.3cm]
{\large\color{isls-gray} CLI-Generierung, Lokales LLM, Selbst-Scraping}\\[2cm]

{\large
\begin{tabular}{rl}
\textbf{System:}    & ISLS --- Intelligent Semantic Ledger Substrate \\
\textbf{Version:}   & v4.2 (D8) \\
\textbf{Author:}    & Sebastian Klemm \\
\textbf{Date:}      & \today \\
\textbf{Status:}    & Specification \\
\textbf{Depends:}   & D7 (complete), InfraBlueprint (D6) \\
\textbf{Branch:}    & \texttt{claude/implement-isls-d8-spec} \\
\end{tabular}
}

\vfill

{\small\color{isls-gray}
Mithras Substructure University\\
Repository: \texttt{https://github.com/LashSesh/graphity}\\
License: MIT
}
\end{titlepage}

\tableofcontents
\newpage


% ══════════════════════════════════════════════════════════════════════════════
\section{Executive Summary}
% ══════════════════════════════════════════════════════════════════════════════

D6 made the HDAG norm-driven. D7 built the cockpit. But ISLS still 
generates exactly one shape: Actix-Web/PostgreSQL/Docker CRUD web 
apps. The \texttt{InfraBlueprint} has flags for CLI, Library, and 
Desktop --- but the structural generators behind those flags do not 
exist. The flags are wires without bulbs.

D8 wires the bulbs. Three concrete expansions:

\begin{enumerate}[nosep]
\item \textbf{CLI Tool Generation:} \texttt{isls forge-chat -m 
      "CLI tool for converting CSV to JSON"} produces a working 
      Rust binary with Clap argument parsing --- no web server, no 
      database, no frontend.
\item \textbf{Ollama Oracle:} Local LLM backend via Ollama HTTP API.
      Zero external API cost. Runs on the Samsung Galaxy Book4 
      with a 7B model.
\item \textbf{Self-Scraping:} \texttt{isls scrape --manifest 
      repos\_self.toml} scrapes Mithra's own 82 repositories, 
      feeding the norm system with patterns from computational 
      engines, trading systems, and parser engines.
\end{enumerate}

These three are not independent features. They form a closed loop: 
the Ollama oracle makes generation free (no API budget). Self-scraping 
feeds the norms. CLI generation proves the system is not limited to 
one shape. Together, they open the horizon.

\subsection{Why These Three}

\begin{itemize}[nosep]
\item \textbf{Ollama:} The OpenAI budget (\textasciitilde\$19) expires 
      May 1, 2026. Without a local LLM backend, ISLS cannot generate 
      after that date. Ollama is the survival condition.
\item \textbf{CLI Generation:} Mithra builds CLI tools (82 repos). 
      The system should generate what the operator actually needs, 
      not just what was easiest to implement first.
\item \textbf{Self-Scraping:} The 82 repos contain the topology of 
      computational engines, trading systems, and cognitive 
      architectures. This is the training corpus no one else has. 
      Feeding it to the norm system is the strategic advantage.
\end{itemize}


% ══════════════════════════════════════════════════════════════════════════════
\section{Work Packages}
% ══════════════════════════════════════════════════════════════════════════════

% ──────────────────────────────────────────────────────────────────────────────
\subsection{W1: CLI Structural Generators}
\label{sec:w1}
% ──────────────────────────────────────────────────────────────────────────────

\subsubsection{Goal}

When the \texttt{InfraBlueprint} has \texttt{has\_cli = true} and 
\texttt{has\_http\_server = false}, the structural generators produce 
a Clap-based CLI binary instead of an Actix-Web server.

\subsubsection{What Changes}

\paragraph{\texttt{generate\_cargo\_toml(spec, bp)}}

Currently always emits actix-web, sqlx, jsonwebtoken, etc. After D8:

\begin{itemize}[nosep]
\item If \texttt{bp.has\_http\_server}: current behavior (actix-web 
      stack).
\item If \texttt{bp.has\_cli \&\& !bp.has\_http\_server}: emit Clap 
      + serde + serde\_json + thiserror + tracing. No actix-web, 
      no sqlx, no bcrypt.
\item If \texttt{bp.has\_database}: add sqlx to CLI deps.
\item Always: serde, tracing, thiserror, chrono.
\end{itemize}

\begin{lstlisting}[style=toml]
# CLI-only Cargo.toml (generated)
[package]
name = "csv-to-json"
version = "0.1.0"
edition = "2021"

[[bin]]
name = "csv-to-json"
path = "src/main.rs"

[dependencies]
clap        = { version = "4", features = ["derive"] }
serde       = { version = "1", features = ["derive"] }
serde_json  = "1"
thiserror   = "1"
tracing     = "0.1"
tracing-subscriber = { version = "0.3", features = ["env-filter"] }
chrono      = { version = "0.4", features = ["serde"] }
\end{lstlisting}

\paragraph{\texttt{generate\_main\_rs(spec, bp)}}

\begin{itemize}[nosep]
\item If \texttt{bp.has\_http\_server}: current behavior 
      (Actix HttpServer).
\item If \texttt{bp.has\_cli \&\& !bp.has\_http\_server}: Clap 
      parser with one subcommand per entity operation.
\end{itemize}

\begin{lstlisting}[style=rust]
// CLI main.rs (generated, structural)
use clap::Parser;

mod errors;
mod models;
mod services;

#[derive(Parser)]
#[command(name = "{app_name}")]
#[command(about = "{description}")]
struct Cli {
    #[command(subcommand)]
    command: Commands,
}

#[derive(clap::Subcommand)]
enum Commands {
    /// Process a CSV file
    Process {
        #[arg(short, long)]
        input: String,
        #[arg(short, long)]
        output: Option<String>,
    },
    /// Show version info
    Version,
}

fn main() {
    tracing_subscriber::fmt::init();
    let cli = Cli::parse();
    
    match cli.command {
        Commands::Process { input, output } => {
            // Entity-driven processing
            let out = output.unwrap_or_else(|| "output.json".into());
            if let Err(e) = services::process(&input, &out) {
                eprintln!("Error: {}", e);
                std::process::exit(1);
            }
        }
        Commands::Version => {
            println!("{} v0.1.0", "{app_name}");
        }
    }
}
\end{lstlisting}

\paragraph{CLI HDAG Structure}

When blueprint is CLI-only, the HDAG omits:
\begin{itemize}[nosep]
\item All \texttt{api/} nodes (no HTTP handlers).
\item All \texttt{frontend/} nodes (no UI).
\item \texttt{auth.rs}, \texttt{pagination.rs} (no web context).
\item \texttt{Dockerfile}, \texttt{docker-compose.yml}, 
      \texttt{nginx.conf} (no deployment).
\item \texttt{database/pool.rs} (unless \texttt{has\_database}).
\end{itemize}

Retains:
\begin{itemize}[nosep]
\item \texttt{Cargo.toml} (with Clap deps).
\item \texttt{main.rs} (with Clap parser).
\item \texttt{src/errors.rs} (error types).
\item \texttt{src/models/*.rs} (data structs).
\item \texttt{src/services/*.rs} (business logic).
\item \texttt{tests/} (unit tests).
\item \texttt{.gitignore}.
\end{itemize}

The LLM files for a CLI tool are the \texttt{services/*.rs} files 
(business logic implementation) instead of \texttt{database/*\_queries.rs} 
(SQL queries). This is a key shift: the LLM surface moves from 
``write SQL'' to ``write business logic.''

\subsubsection{Modified Files}

\begin{longtable}{lp{10cm}}
\toprule
\textbf{File} & \textbf{Change} \\
\midrule
\texttt{isls-forge-llm/src/static\_files.rs}
  & \texttt{generate\_cargo\_toml(spec, bp)} --- CLI branch with 
    Clap deps. \\
\texttt{isls-forge-llm/src/structural.rs}
  & \texttt{generate\_main\_rs(spec, bp)} --- Clap parser. 
    \texttt{generate\_cli\_service\_rs(entity)} --- placeholder 
    service for CLI entities. \\
\texttt{isls-forge-llm/src/hdag.rs}
  & \texttt{build(spec, bp)} --- CLI branch omits web nodes, 
    makes services LLM-generated instead of structural. \\
\bottomrule
\end{longtable}

\subsubsection{W1 Verification}

\begin{enumerate}[nosep]
\item \texttt{isls forge-chat -m "CLI tool for converting CSV 
      files to JSON" --mock-oracle --output ./output/csv-tool}
\item Verify output has: \texttt{Cargo.toml} with Clap, 
      \texttt{main.rs} with Parser, NO \texttt{api/}, NO 
      \texttt{frontend/}, NO \texttt{Dockerfile}.
\item \texttt{cd output/csv-tool \&\& cargo build} --- compiles.
\item \texttt{isls forge-chat -m "Restaurant app with menus"} --- 
      still generates web app (backward compat).
\end{enumerate}


% ──────────────────────────────────────────────────────────────────────────────
\subsection{W2: Ollama Oracle}
\label{sec:w2}
% ──────────────────────────────────────────────────────────────────────────────

\subsubsection{Goal}

Add an \texttt{OllamaOracle} implementation of the \texttt{Oracle} 
trait that calls a locally running Ollama instance via HTTP API.

\subsubsection{Ollama API}

Ollama exposes a simple HTTP API at \texttt{http://localhost:11434}:

\begin{lstlisting}[style=bash]
POST http://localhost:11434/api/generate
Content-Type: application/json

{
  "model": "codellama:7b",
  "prompt": "Generate Rust code for...",
  "stream": false,
  "options": {
    "temperature": 0.1,
    "num_predict": 4096
  }
}
\end{lstlisting}

Response:
\begin{lstlisting}[style=bash]
{
  "response": "use sqlx::PgPool;\n...",
  "done": true
}
\end{lstlisting}

\subsubsection{Implementation}

\begin{lstlisting}[style=rust]
pub struct OllamaOracle {
    pub model: String,
    pub base_url: String,
    client: reqwest::blocking::Client,
}

impl OllamaOracle {
    pub fn new(model: &str, base_url: &str) -> Self {
        Self {
            model: model.to_string(),
            base_url: base_url.to_string(),
            client: reqwest::blocking::Client::builder()
                .timeout(std::time::Duration::from_secs(300))
                .build()
                .expect("HTTP client"),
        }
    }
}

impl Oracle for OllamaOracle {
    fn call(&self, prompt: &str, max_tokens: u32) 
        -> Result<String, Box<dyn std::error::Error>> 
    {
        let url = format!("{}/api/generate", self.base_url);
        let body = serde_json::json!({
            "model": self.model,
            "prompt": prompt,
            "stream": false,
            "options": {
                "temperature": 0.1,
                "num_predict": max_tokens
            }
        });
        let resp = self.client.post(&url)
            .json(&body)
            .send()?;
        let json: serde_json::Value = resp.json()?;
        let response = json["response"].as_str()
            .ok_or("no response field")?;
        Ok(response.to_string())
    }
}
\end{lstlisting}

\subsubsection{CLI Integration}

New CLI option across all forge commands:

\begin{lstlisting}[style=bash]
isls forge-chat -m "..." --ollama --model codellama:7b
isls forge-v2 --requirements ... --ollama
isls forge-self --ollama --model mistral:7b
\end{lstlisting}

\texttt{--ollama} activates the Ollama oracle. Default model: 
\texttt{codellama:7b}. Default URL: \texttt{http://localhost:11434}.
Overridable via \texttt{--ollama-url}.

Oracle selection priority:
\begin{enumerate}[nosep]
\item \texttt{--mock-oracle} $\to$ MockOracle (no LLM).
\item \texttt{--ollama} $\to$ OllamaOracle (local).
\item \texttt{--api-key} or \texttt{OPENAI\_API\_KEY} $\to$ 
      OpenAiOracle (remote).
\item None $\to$ MockOracle (fallback).
\end{enumerate}

\subsubsection{Availability Check}

Before using Ollama, verify it is running:

\begin{lstlisting}[style=bash]
GET http://localhost:11434/api/tags
\end{lstlisting}

If the request fails: print ``Ollama is not running. Start it 
with `ollama serve' or install from https://ollama.com'' and fall 
back to MockOracle.

If the requested model is not pulled: print ``Model \{model\} not 
found. Pull it with `ollama pull \{model\}''' and fail.

\subsubsection{New File}

\begin{longtable}{lp{10cm}}
\toprule
\textbf{File} & \textbf{Purpose} \\
\midrule
\texttt{isls-forge-llm/src/oracle/ollama.rs}
  & \texttt{OllamaOracle} struct + \texttt{Oracle} impl. \\
\bottomrule
\end{longtable}

\subsubsection{Modified Files}

\begin{longtable}{lp{10cm}}
\toprule
\textbf{File} & \textbf{Change} \\
\midrule
\texttt{isls-forge-llm/src/oracle.rs} (or \texttt{oracle/mod.rs})
  & Re-export OllamaOracle. \\
\texttt{isls-cli/src/main.rs}
  & Add \texttt{--ollama}, \texttt{--ollama-url}, 
    \texttt{--model} to all forge commands. Oracle selection 
    logic. \\
\bottomrule
\end{longtable}

\subsubsection{W2 Verification}

\begin{enumerate}[nosep]
\item Install Ollama, pull \texttt{codellama:7b}.
\item \texttt{isls forge-chat -m "Pet shop app" --ollama 
      --output ./output/ollama-test}
\item Verify LLM calls go to localhost:11434 (check Ollama logs).
\item Verify generated app compiles (or at least S6 Coagula 
      attempts fixes).
\item \texttt{isls forge-chat -m "Pet shop app" --api-key \$KEY}
      --- still works with OpenAI (backward compat).
\item Without Ollama running: verify graceful error message.
\end{enumerate}


% ──────────────────────────────────────────────────────────────────────────────
\subsection{W3: Self-Scraping Manifest}
\label{sec:w3}
% ──────────────────────────────────────────────────────────────────────────────

\subsubsection{Goal}

Create a scrape manifest listing Mithra's own repositories and 
run it to seed the norm system with patterns from computational 
engines, trading systems, and cognitive architectures.

\subsubsection{Manifest}

File: \texttt{examples/scrape\_self.toml}

\begin{lstlisting}[style=toml]
# ISLS Self-Scraping Manifest
# Mithra's own repositories — the corpus no one else has.

[[repo]]
path = "../fsr-engine"
domain = "fsr-trading"

[[repo]]
path = "../babylon-compiler"
domain = "babylon-cuneiform"

[[repo]]
path = "../tablet-forge"
domain = "tablet-forge"

[[repo]]
path = "../pse-core"
domain = "pse-parser"

[[repo]]
url = "https://github.com/LashSesh/graphity.git"
domain = "graphity-isls"
\end{lstlisting}

The manifest uses relative paths (assuming repos are siblings of 
graphity). Missing repos are skipped gracefully (D5 already handles 
this).

\subsubsection{Execution}

\begin{lstlisting}[style=bash]
isls scrape --manifest examples/scrape_self.toml
\end{lstlisting}

This scrapes each local repo, extracts topology, feeds 
\texttt{observe\_and\_learn()}, and reports what was found.

\subsubsection{Expected Norm Impact}

Computational engine repos have different topology than CRUD apps:
\begin{itemize}[nosep]
\item \texttt{types/} instead of \texttt{models/}
\item \texttt{operators/} instead of \texttt{services/}
\item \texttt{solvers/} instead of \texttt{api/}
\item No \texttt{frontend/}, no \texttt{migrations/}
\end{itemize}

This will create new cross-layer patterns (e.g. 
Types+Operators+Solvers) that are distinct from the CRUD patterns 
(Model+Query+Service+Api+Frontend). Over time, these patterns 
promote to auto-norms that describe computational engine topology.

\subsubsection{W3 Verification}

\begin{enumerate}[nosep]
\item \texttt{isls scrape --manifest examples/scrape\_self.toml}
\item Verify at least 2--3 repos scraped successfully.
\item \texttt{isls norms candidates} --- new candidates from 
      non-CRUD topologies.
\item \texttt{isls norms stats} --- observation count increased, 
      domain list includes engine domains.
\end{enumerate}


% ──────────────────────────────────────────────────────────────────────────────
\subsection{W4: Pipeline + Studio Bat Fix}
\label{sec:w4}
% ──────────────────────────────────────────────────────────────────────────────

\subsubsection{Goal}

Fix the \texttt{isls\_studio.bat} API key handling so the key is 
correctly passed to the server, and extend the pipeline bat with 
a CLI generation test and an Ollama test (if available).

\subsubsection{Studio Bat Fix}

The \texttt{isls\_studio.bat} currently does not pass the API key 
correctly to the \texttt{isls serve} command. Fix:

\begin{lstlisting}[style=bash]
rem After API key prompt:
set OPENAI_API_KEY=%API_KEY%
target\release\isls.exe serve --port 8420 --api-key %API_KEY%
\end{lstlisting}

Ensure the API key is both set as environment variable AND passed 
as CLI argument (the server checks both).

\subsubsection{Pipeline Extension}

Add to \texttt{isls\_pipeline.bat} after D6 steps:

\begin{lstlisting}[style=bash]
rem ── D8: CLI Generation Test ──
isls forge-chat -m "CLI tool that counts words in text files" ^
    --api-key %API_KEY% --output ./output/pipeline_%TIMESTAMP%/cli-wordcount
cd output\pipeline_%TIMESTAMP%\cli-wordcount && cargo build
rem Result: PASS if compiles, FAIL with log

rem ── D8: Ollama Test (optional) ──
rem Check if Ollama is running
curl -s http://localhost:11434/api/tags >nul 2>&1
if %ERRORLEVEL%==0 (
    isls forge-chat -m "Simple hello world app" --ollama ^
        --output ./output/pipeline_%TIMESTAMP%/ollama-test
) else (
    echo Ollama not running - skipping test
)
\end{lstlisting}

\subsubsection{W4 Verification}

\begin{enumerate}[nosep]
\item \texttt{isls\_studio.bat} --- API key is prompted and 
      correctly used by the server.
\item \texttt{isls\_pipeline.bat} --- includes CLI generation 
      test in the report.
\end{enumerate}


% ══════════════════════════════════════════════════════════════════════════════
\section{File Inventory}
% ══════════════════════════════════════════════════════════════════════════════

\subsection{New Files}

\begin{longtable}{llp{7cm}}
\toprule
\textbf{Crate} & \textbf{File} & \textbf{Purpose} \\
\midrule
isls-forge-llm & \texttt{src/oracle/ollama.rs}
  & Ollama Oracle implementation \\
isls & \texttt{examples/scrape\_self.toml}
  & Self-scraping manifest \\
\bottomrule
\end{longtable}

\subsection{Modified Files}

\begin{longtable}{llp{7cm}}
\toprule
\textbf{Crate} & \textbf{File} & \textbf{Change} \\
\midrule
isls-forge-llm & \texttt{src/static\_files.rs}
  & Blueprint-aware Cargo.toml (CLI branch) \\
isls-forge-llm & \texttt{src/structural.rs}
  & CLI main.rs with Clap, CLI service stubs \\
isls-forge-llm & \texttt{src/hdag.rs}
  & CLI branch in build() --- no web nodes, services as LLM \\
isls-forge-llm & \texttt{src/oracle.rs}
  & Re-export OllamaOracle \\
isls-cli & \texttt{src/main.rs}
  & --ollama flag, oracle selection, CLI test in pipeline \\
(root) & \texttt{isls\_pipeline.bat}
  & CLI generation test, Ollama test \\
(root) & \texttt{isls\_studio.bat}
  & API key fix \\
\bottomrule
\end{longtable}


% ══════════════════════════════════════════════════════════════════════════════
\section{The 12 Rules --- D8 Compliance}
% ══════════════════════════════════════════════════════════════════════════════

\begin{longtable}{clp{9cm}}
\toprule
\textbf{\#} & \textbf{Status} & \textbf{D8 Impact} \\
\midrule
1  & \color{isls-green}\checkmark & CLI mode shifts LLM surface 
     from queries to services --- still 7 or fewer LLM files. \\
2  & \color{isls-green}\checkmark & mod.rs generated for populated 
     directories. \\
3  & \color{isls-green}\checkmark & Entity naming convention 
     preserved. \\
4  & \color{isls-green}\checkmark & No auth in CLI-only mode. \\
5--6 & \color{isls-green}\checkmark & No query files in CLI-only 
     mode (unless has\_database). \\
7  & \color{isls-green}\checkmark & No bcrypt in CLI-only deps. 
     Tracing always present. \\
8  & \color{isls-green}\checkmark & No seed in CLI-only mode. \\
9  & \color{isls-green}\checkmark & No Dockerfile in CLI-only mode.\\
10 & \color{isls-green}\checkmark & Ollama unavailable $\to$ 
     graceful error + fallback. \\
11 & \color{isls-green}\checkmark & ProvidedSymbols adapt to 
     CLI context (no PaginationParams, no AuthUser). \\
12 & \color{isls-green}\checkmark & LLM prompts include full 
     type context regardless of mode. \\
\bottomrule
\end{longtable}


% ══════════════════════════════════════════════════════════════════════════════
\section{Constraints for Claude Code}
% ══════════════════════════════════════════════════════════════════════════════

\begin{enumerate}
\item \textbf{Backward compatibility.} All D1--D7 apps MUST still 
      generate identically. Test with the pipeline bat.
\item \textbf{Ollama uses \texttt{reqwest::blocking}}, which is 
      already in the workspace. No new HTTP crate.
\item \textbf{The oracle module} may need restructuring from a 
      single file to a module directory 
      (\texttt{oracle/mod.rs}, \texttt{oracle/openai.rs}, 
      \texttt{oracle/ollama.rs}, \texttt{oracle/mock.rs}).
      This is acceptable as long as the public API 
      (\texttt{Oracle} trait, \texttt{MockOracle}, 
      \texttt{OpenAiOracle}) stays unchanged.
\item \textbf{CLI structural generators} must be conditional on 
      blueprint flags. When \texttt{has\_cli \&\& !has\_http\_server}, 
      produce CLI shape. When \texttt{has\_http\_server}, produce 
      web shape (existing behavior).
\item \textbf{Do not modify} isls-norms, isls-reader, isls-code-topo, 
      isls-hypercube, isls-types, isls-agent, isls-gateway.
\item \textbf{The self-scraping manifest} must use relative paths 
      and handle missing repos gracefully.
\item All new code MUST compile with \texttt{cargo build --workspace}.
\item \texttt{cargo test --workspace} --- all existing tests pass 
      (218+) plus new ones.
\item Branch: \texttt{claude/implement-isls-d8-spec}.
\item Commit format: \texttt{D8/W\{n\}: <description>}.
\end{enumerate}


% ══════════════════════════════════════════════════════════════════════════════
\section{Dependency Graph}
% ══════════════════════════════════════════════════════════════════════════════

\begin{verbatim}
W1 (CLI Structural Generators)     W2 (Ollama Oracle)
         |                                |
         v                                v
W4 (Pipeline + Bat Fixes) <──────────────┘
         |
         v
W3 (Self-Scraping Manifest)
\end{verbatim}

W1 and W2 are independent. W4 integrates both into the pipeline.
W3 runs last (uses the completed system).

Recommended: W2 $\to$ W1 $\to$ W4 $\to$ W3.
(Ollama first --- eliminates API cost for all subsequent testing.)


% ══════════════════════════════════════════════════════════════════════════════
\section{Acceptance Criteria}
% ══════════════════════════════════════════════════════════════════════════════

\begin{longtable}{clp{9cm}}
\toprule
\textbf{\#} & \textbf{Pass} & \textbf{Criterion} \\
\midrule
1 & & \texttt{cargo build --workspace} compiles. \\
2 & & \texttt{cargo test --workspace} passes (218+). \\
3 & & CLI generation: \texttt{forge-chat -m "CLI tool..."} 
      produces Clap binary, compiles. \\
4 & & Web app generation still works (6 D4 domains compile). \\
5 & & Ollama oracle: \texttt{forge-chat --ollama} calls 
      localhost:11434. \\
6 & & Without Ollama: graceful error, no crash. \\
7 & & Self-scraping: \texttt{scrape --manifest 
      examples/scrape\_self.toml} runs. \\
8 & & \texttt{isls\_studio.bat} --- API key correctly passed. \\
9 & & Pipeline bat includes CLI generation test. \\
\bottomrule
\end{longtable}

D8 is passed when criteria 1--4 and 6 are fulfilled.
Criteria 5 requires Ollama installed (manual test).
Criteria 7 requires local repos (manual test).


% ══════════════════════════════════════════════════════════════════════════════
\section{Post-D8: What Opens}
% ══════════════════════════════════════════════════════════════════════════════

With D8 complete, the Donnerkeil roadmap is finished. But the system 
is not. What opens:

\begin{description}[nosep]
\item[Library Generation] \texttt{has\_library = true} --- lib.rs 
  with pub API, no binary. Generates reusable crates.
\item[Desktop Generation] \texttt{has\_gui = true} --- egui or iced 
  desktop app. No web server. Local database.
\item[Multi-Infrastructure] \texttt{has\_http\_server + has\_cli} ---
  server with CLI management. Both in one binary.
\item[Norm Depth] More norms, finer granularity. Connection pooling, 
  caching strategies, error boundaries, observability patterns.
\item[Self-Improvement] The Cockpit (D7) generates apps. Those apps 
  are observed (D4). Observations feed norms. Norms improve future 
  generations. The spiral tightens autonomously.
\end{description}

The 8 Donnerkeil-Durchschl\"age built the machine. Everything after 
is the machine using itself.

\vfill
\begin{center}
\color{isls-gray}\rule{0.5\textwidth}{0.4pt}\\[0.5em]
{\small End of D8 Specification.\\[0.3em]
Der Horizont ist offen.\\
Der Spielautomat hat seine Walzen.\\
Die erste Walze, die nicht CRUD ist, dreht sich.}
\end{center}

\end{document}
